\documentclass{amsart}
% For dealing with references we use the comment environment
\usepackage{verbatim}
\newenvironment{reference}{\comment}{\endcomment}
%\newenvironment{reference}{}{}
\newenvironment{slogan}{\comment}{\endcomment}
\newenvironment{history}{\comment}{\endcomment}

% For commutative diagrams we use Xy-pic
\usepackage[all]{xy}

% We use 2cell for 2-commutative diagrams.
\xyoption{2cell}
\UseAllTwocells

% We use multicol for the list of chapters between chapters
\usepackage{multicol}

% This is generally recommended for better output
\usepackage{lmodern}
\usepackage[T1]{fontenc}

% For cross-file-references

% Package for hypertext links:
\usepackage{hyperref}

% For any local file, say "hello.tex" you want to link to please

% Theorem environments.
%
\theoremstyle{plain}
\newtheorem{theorem}[subsection]{Theorem}
\newtheorem{proposition}[subsection]{Proposition}
\newtheorem{lemma}[subsection]{Lemma}

\theoremstyle{definition}
\newtheorem{definition}[subsection]{Definition}
\newtheorem{example}[subsection]{Example}
\newtheorem{exercise}[subsection]{Exercise}
\newtheorem{situation}[subsection]{Situation}

\theoremstyle{remark}
\newtheorem{remark}[subsection]{Remark}
\newtheorem{remarks}[subsection]{Remarks}

\numberwithin{equation}{subsection}

% Macros
%
\def\lim{\mathop{\mathrm{lim}}\nolimits}
\def\colim{\mathop{\mathrm{colim}}\nolimits}
\def\Spec{\mathop{\mathrm{Spec}}}
\def\Hom{\mathop{\mathrm{Hom}}\nolimits}
\def\Ext{\mathop{\mathrm{Ext}}\nolimits}
\def\SheafHom{\mathop{\mathcal{H}\!\mathit{om}}\nolimits}
\def\SheafExt{\mathop{\mathcal{E}\!\mathit{xt}}\nolimits}
\def\Sch{\mathit{Sch}}
\def\Mor{\mathop{\mathrm{Mor}}\nolimits}
\def\Ob{\mathop{\mathrm{Ob}}\nolimits}
\def\Sh{\mathop{\mathit{Sh}}\nolimits}
\def\NL{\mathop{N\!L}\nolimits}
\def\CH{\mathop{\mathrm{CH}}\nolimits}
\def\proetale{{pro\text{-}\acute{e}tale}}
\def\etale{{\acute{e}tale}}
\def\QCoh{\mathit{QCoh}}
\def\Ker{\mathop{\mathrm{Ker}}}
\def\Im{\mathop{\mathrm{Im}}}
\def\Coker{\mathop{\mathrm{Coker}}}
\def\Coim{\mathop{\mathrm{Coim}}}

% Boxtimes
%
\DeclareMathSymbol{\boxtimes}{\mathbin}{AMSa}{"02}

%
% Macros for moduli stacks/spaces
%
\def\QCohstack{\mathcal{QC}\!\mathit{oh}}
\def\Cohstack{\mathcal{C}\!\mathit{oh}}
\def\Spacesstack{\mathcal{S}\!\mathit{paces}}
\def\Quotfunctor{\mathrm{Quot}}
\def\Hilbfunctor{\mathrm{Hilb}}
\def\Curvesstack{\mathcal{C}\!\mathit{urves}}
\def\Polarizedstack{\mathcal{P}\!\mathit{olarized}}
\def\Complexesstack{\mathcal{C}\!\mathit{omplexes}}
% \Pic is the operator that assigns to X its picard group, usage \Pic(X)
% \Picardstack_{X/B} denotes the Picard stack of X over B
% \Picardfunctor_{X/B} denotes the Picard functor of X over B
\def\Pic{\mathop{\mathrm{Pic}}\nolimits}
\def\Picardstack{\mathcal{P}\!\mathit{ic}}
\def\Picardfunctor{\mathrm{Pic}}
\def\Deformationcategory{\mathcal{D}\!\mathit{ef}}

\setlength{\emergencystretch}{3em}
\begin{document}
\title[Stacks Project, Tag 0FBT]{457 --- Smooth targets identify bivariant Chow groups with degree-shifted bivariant groups}
\maketitle
\noindent Copyright (C) 2005--2025 Johan de Jong. Stacks Project authors. Source: \href{https://stacks.math.columbia.edu/tag/0FBT}{Tag 0FBT}.

\noindent Editorial difficulty note (not author text). Difficulty H2: The inverse to composition with a smooth Gysin class must be constructed rather than inferred from flat pullback. The proof uses the relative diagonal and two cartesian restriction diagrams, then verifies both compositions by diagonal-section identities and bivariant Gysin compatibility..

\noindent Editorial provenance note (not author text). History: excerpt from revision \texttt{a0444\allowbreak 6e57e\allowbreak c1fbc\allowbreak 252a8\allowbreak 71afc\allowbreak ec775\allowbreak 2fb28\allowbreak 07b14}, chow.tex, lines 14002--14095. Reference presentation was adapted; original-author context and core support blocks were retained or added. The mathematical wording is unchanged. Licensed under GNU FDL 1.2 or later; no invariant sections or cover texts. See ../licenses/COPYING. Context blocks reproduce author definitions and supporting results; they are not additional counted problems.

\begin{definition}
\label{definition-bivariant-class}
\begin{reference}
Similar to \href{https://stacks.math.columbia.edu/bibliography}{Bibliography: F (Definition 17.1)}
\end{reference}
Let $(S, \delta)$ be as in Situation \ref{situation-setup}.
Let $f : X \to Y$ be a morphism of schemes locally of finite type over $S$.
Let $p \in \mathbf{Z}$.
A {\it bivariant class $c$ of degree $p$ for $f$} is given by a rule
which assigns to every locally of finite type morphism $Y' \to Y$
and every $k$ a map
$$
c \cap - : \CH_k(Y') \longrightarrow \CH_{k - p}(X')
$$
where $X' = Y' \times_Y X$, satisfying the following conditions
\begin{enumerate}
\item if $Y'' \to Y'$ is a proper, then
$c \cap (Y'' \to Y')_*\alpha'' = (X'' \to X')_*(c \cap \alpha'')$
for all $\alpha''$ on $Y''$ where $X'' = Y'' \times_Y X$,
\item if $Y'' \to Y'$ is flat locally of finite type of
fixed relative dimension, then
$c \cap (Y'' \to Y')^*\alpha' = (X'' \to X')^*(c \cap \alpha')$
for all $\alpha'$ on $Y'$, and
\item if $(\mathcal{L}', s', i' : D' \to Y')$ is as in
Definition \ref{definition-gysin-homomorphism}
with pullback $(\mathcal{N}', t', j' : E' \to X')$ to $X'$,
then we have $c \cap (i')^*\alpha' = (j')^*(c \cap \alpha')$
for all $\alpha'$ on $Y'$.
\end{enumerate}
The collection of all bivariant classes of degree $p$ for $f$ is
denoted $A^p(X \to Y)$.
\end{definition}

\begin{situation}
\label{situation-setup}
Here $S$ is a locally Noetherian, and universally catenary scheme.
Moreover, we assume $S$ is endowed with a dimension function
$\delta : S \longrightarrow \mathbf{Z}$.
\end{situation}

\begin{definition}
\label{definition-gysin-homomorphism}
Let $(S, \delta)$ be as in Situation \ref{situation-setup}.
Let $X$ be locally of finite type over $S$.
Let $(\mathcal{L}, s)$ be a pair consisting of an invertible
sheaf and a global section $s \in \Gamma(X, \mathcal{L})$.
Let $D = Z(s)$ be the zero scheme of $s$, and
denote $i : D \to X$ the closed immersion.
We define, for every integer $k$, a {\it Gysin homomorphism}
$$
i^* : Z_{k + 1}(X) \to \CH_k(D).
$$
by the following rules:
\begin{enumerate}
\item Given an integral closed subscheme $W \subset X$ with
$\dim_\delta(W) = k + 1$ we define
\begin{enumerate}
\item if $W \not \subset D$, then $i^*[W] = [D \cap W]_k$ as a
$k$-cycle on $D$, and
\item if $W \subset D$, then
$i^*[W] = i'_*(c_1(\mathcal{L}|_W) \cap [W])$,
where $i' : W \to D$ is the induced closed immersion.
\end{enumerate}
\item For a general $(k + 1)$-cycle $\alpha = \sum n_j[W_j]$
we set
$$
i^*\alpha = \sum n_j i^*[W_j]
$$
\item If $D$ is an effective Cartier divisor, then we denote
$D \cdot \alpha = i_*i^*\alpha$ the pushforward of the class $i^*\alpha$
to a class on $X$.
\end{enumerate}
\end{definition}

\begin{proposition}
\label{proposition-compute-bivariant}
\begin{reference}
\href{https://stacks.math.columbia.edu/bibliography}{Bibliography: F (Proposition 17.4.2)}
\end{reference}
Let $(S, \delta)$ be as in Situation \ref{situation-setup}.
Let $f : X \to Y$ and $g : Y \to Z$ be morphisms of schemes locally
of finite type over $S$. If $g$ is smooth of relative dimension $d$, then
$A^p(X \to Y) = A^{p - d}(X \to Z)$.
\end{proposition}

\begin{proof}
We will use that smooth morphisms are local complete intersection
morphisms whose Gysin maps exist (see Section \href{https://stacks.math.columbia.edu/tag/0FEZ}{section koszul (Tag 0FEZ)}).
In particular we have $g^! \in A^{-d}(Y \to Z)$. Then we can send
$c \in A^p(X \to Y)$ to $c \circ g^! \in A^{p - d}(X \to Z)$.

\medskip\noindent
Conversely, let $c' \in A^{p - d}(X \to Z)$. Denote $res(c')$ the restriction
(Remark \href{https://stacks.math.columbia.edu/tag/0F9Z}{remark restriction bivariant (Tag 0F9Z)}) of $c'$ by the morphism $Y \to Z$.
Since the diagram
$$
\xymatrix{
X \times_Z Y \ar[r]_{\text{pr}_2} \ar[d]_{\text{pr}_1} & Y \ar[d]^g \\
X \ar[r]^f & Z
}
$$
is cartesian we find $res(c') \in A^{p - d}(X \times_Z Y \to Y)$.
Let $\Delta : Y \to Y \times_Z Y$ be the diagonal and denote
$res(\Delta^!)$ the restriction of $\Delta^!$
to $X \times_Z Y$ by the morphism $X \times_Z Y \to Y \times_Z Y$.
Since the diagram
$$
\xymatrix{
X \ar[r] \ar[d] & X \times_Z Y \ar[d] \\
Y \ar[r]^-\Delta & Y \times_Z Y
}
$$
is cartesian we see that $res(\Delta^!) \in A^d(X \to X \times_Z Y)$.
Combining these two restrictions we obtain
$$
res(\Delta^!) \circ res(c') \in A^p(X \to Y)
$$
Thus we have produced maps $A^p(X \to Y) \to A^{p - d}(X \to Z)$
and $A^{p - d}(X \to Z) \to A^p(X \to Y)$. To finish the proof we
will show these maps are mutually inverse.

\medskip\noindent
Let us start with $c \in A^p(X \to Y)$. Consider the diagram
$$
\xymatrix{
X \ar[d] \ar[r] & Y \ar[d] \\
X \times_Z Y \ar[r] \ar[d]^{\text{pr}_1} &
Y \times_Z Y \ar[r]_{p_2} \ar[d]^{p_1} &
Y \ar[d]^g \\
X \ar[r]^f &
Y \ar[r]^g &
Z
}
$$
whose squares are carteisan. The lower two square of this diagram
show that $res(c \circ g^!) = res(c) \cap p_2^!$ where in this formula
$res(c)$ means the restriction of $c$ via $p_1$. Looking at the upper
square of the diagram and using Lemma \href{https://stacks.math.columbia.edu/tag/0FF6}{lemma lci gysin commutes (Tag 0FF6)}
we get $c \circ \Delta^! = res(\Delta^!) \circ res(c)$.
We compute
\begin{align*}
res(\Delta^!) \circ res(c \circ g^!)
& =
res(\Delta^!) \circ res(c) \circ p_2^! \\
& =
c \circ \Delta^! \circ p_2^! \\
& =
c
\end{align*}
The final equality by Lemma \href{https://stacks.math.columbia.edu/tag/0FBS}{lemma diagonal identity (Tag 0FBS)}.

\medskip\noindent
Conversely, let us start with $c' \in A^{p - d}(X \to Z)$. Looking
at the lower rectangle of the diagram above we find
$res(c') \circ g^! = \text{pr}_1^! \circ c'$.
We compute
\begin{align*}
res(\Delta^!) \circ res(c') \circ g^!
& =
res(\Delta^!) \circ \text{pr}_1^! \circ c' \\
& =
c'
\end{align*}
The final equality holds because the left two squares of
the diagram show that
$\text{id} = res(\Delta^! \circ p_1^!) = res(\Delta^!) \circ \text{pr}_1^!$.
This finishes the proof.
\end{proof}
\end{document}
